\section{The intermediate language and program extraction}
\label{sec:lambdabox}

The first half of the translation is Letouzey-style program extraction:
an erasure of $\CICi$ terms into an untyped functional calculus $\lbox$.
The calculus is \emph{the same} as in the companion paper --- indexed
families change constructor \emph{arities} (a constructor keeps only its
informative arguments) but not the shape of the target language --- and
the erasure function extends that of \cite{Extraction2026} with clauses
for the new formers.  Two design points carry over: since the target of
the overall translation is an equational theory rather than a programming
language with an evaluation order, erased arguments are \emph{dropped}
rather than replaced by dummies; and the singleton eliminations collapse
to their branches structurally.  A third point is new: erasure is
\emph{index-blind}.  Neither the motive indices of a case, nor the index
equations of its branches, nor the index terms of types leave any trace;
index terms occurring as \emph{constructor arguments} (the length
argument of $\csvcons$) are informative and are kept, exactly as in
Letouzey's extraction \cite{Letouzey2003,Letouzey2008}, where redundant
stored indices are likewise kept.  The re-materialisation of index
information is the job of specification extraction
(\cref{sec:translation-spec}), not of the program.

\subsection{The calculus \texorpdfstring{$\lbox$}{lambda-box}}

\begin{definition}[Syntax of $\lbox$]
\label{def:lbox-syntax}
Fix the datatype constructors $\wcns{c}$ of a $\CICi$ environment
$\Env$, each with arity $r$ equal to its number of \emph{informative}
arguments (premises do not count), the term-definition constants $c$ of
$\Env$, and a distinguished constant $\bx$ (``dead code'').  The terms of
$\lbox$ are
\[
  e \Coloneqq x \mid c \mid \bx \mid \lambda x.e \mid e\,e'
    \mid \wcns{c}(e_1, \dots, e_r)
    \mid \mathsf{case}(e;\, \{\wcns{c}_i(\tup{x}_i) \Rightarrow e_i\}_i)
    \mid \mathsf{fix}\,f.\,e
\]
where constructors are fully applied and a case expression carries one
branch for each constructor of the datatype of its constructors.
\end{definition}

\begin{definition}[Reduction of $\lbox$]
\label{def:lbox-red}
The relation $\lred$ is the contextual closure of:
\begin{align*}
  (\lambda x.e)\,e' &\lred \subst{e}{e'}{x}
  &&(\beta)\\
  \mathsf{case}(\wcns{c}_i(\tup{e});\, \{\wcns{c}_j(\tup{x}_j) \Rightarrow
    e_j\}_j) &\lred \subst{e_i}{\tup{e}}{\tup{x}_i}
  &&(\iota)\\
  \mathsf{fix}\,f.\,e &\lred \subst{e}{\mathsf{fix}\,f.e}{f}
  &&(\mu)\\
  c &\lred \erase{t} \quad \text{if } (c \defeq t : \sigma) \in \Env
  &&(\delta)
\end{align*}
where $\erasef$ is the erasure of \cref{def:erasure} (the $\delta$-rule
is well defined because the body of $c$ mentions only constants declared
before $c$).  We write $\lconv$ for the least congruence containing
$\lred$.
\end{definition}

The four rule families have pairwise distinct, non-overlapping,
left-linear redex shapes with rigid heads; $\lbox$ is an orthogonal
higher-order rewrite system, and its confluence is independent of the
constructor arities:

\begin{theorem}[Confluence of $\lbox$]
\label{thm:confluence}
$\lred$ is confluent.  Consequently, $e_1 \lconv e_2$ if and only if
$e_1$ and $e_2$ have a common $\lreds$-reduct.
\end{theorem}

\begin{proof}
By the Tait--Martin-L{\"o}f parallel-reduction method in Takahashi's
formulation \cite{Takahashi1995}; the full proof is in
\cref{sec:confluence}.  Alternatively, $\lbox$ is an orthogonal
combinatory reduction system, so confluence follows from Klop's theorem
\cite{Klop1980}.
\end{proof}

\begin{corollary}[Constructor discrimination and injectivity]
\label{cor:discrimination}
Let $\wcns{c}, \wcns{d}$ be constructors.
\begin{enumerate}[leftmargin=2em]
\item If $\wcns{c}(\tup{e}) \lconv \wcns{d}(\tup{e}')$ then $\wcns{c} =
  \wcns{d}$ and $e_i \lconv e'_i$ for all $i$.
\item An application $e_1\,e_2$ in which $e_1$ is normal and not a
  $\lambda$-abstraction (for instance $e_1 = \wcns{c}(\tup{e})$ with
  $\tup{e}$ normal) is not convertible to any constructor form
  $\wcns{d}(\tup{e}')$.
\end{enumerate}
\end{corollary}

\begin{proof}
(1) Constructor-headed terms only reduce internally: every reduct of
$\wcns{c}(\tup{e})$ has the form $\wcns{c}(\tup{e}'')$ with $e_i \lreds
e''_i$.  By \cref{thm:confluence}, the two sides have a common reduct,
which is then simultaneously $\wcns{c}$-headed and $\wcns{d}$-headed,
forcing $\wcns{c} = \wcns{d}$; its arguments are common reducts of the
respective $e_i, e'_i$.
(2) Such an application is $\lred$-normal apart from reductions inside
$e_1$ and $e_2$, which preserve the application shape, so a common
reduct with a constructor form is impossible as in (1).
\end{proof}

\subsection{Erasure}

\begin{definition}[Erasure]
\label{def:erasure}
The function $\erasef$ from $\CICi$ terms to $\lbox$ terms is defined by
structural recursion:
\begin{align*}
  \erase{x} &= x
  & \erase{c} &= c
  \\
  \erase{t\,a} &= \erase{t}\,\erase{a}
  & \erase{t\,\pi} &= \erase{t}
  \\
  \erase{\lambda x{:}\sigma.t} &= \lambda x.\erase{t}
  & \erase{\lambda p{:}\varphi.t} &= \erase{t}
  \\
  \erase{\wcns{c}(\tup{a};\tup{\pi})} &= \wcns{c}(\erase{\tup{a}})
  & \erase{\tfix\,f{:}T.t} &= \mathsf{fix}\,f.\erase{t}
  \\
  \erase{\tcase{a}{\tup{z}.z.\tau}{\{\wcns{c}_i(\tup{y}_i;\tup{p}_i;\tup{q}_i)
    \Rightarrow t_i\}_i}}
    &\mathrlap{{}= \mathsf{case}(\erase{a};\,
       \{\wcns{c}_i(\tup{y}_i) \Rightarrow \erase{t_i}\}_i)}
  \\
  \erase{\trefine{t}{\pi}} &= \erase{t}
  & \erase{\tval{t}} &= \erase{t}
  \\
  \erase{\tsplit{\pi}{p_1 p_2.t}} &= \erase{t}
  & \erase{\texfalso{\pi}{\sigma}} &= \bx
  \\
  \erase{\tcast{\pi}{z.\tau}{t}} &= \erase{t}
  \\
  \erase{\tscase{J}{\pi}{\tup{z}.\tau}{\tup{p}.t}}
    &\mathrlap{{}= \erase{t}}
\end{align*}
\end{definition}

The clauses for the new formers say: constructor terms keep their
informative arguments and lose their premise proofs; a case on a family
keeps its scrutinee and branch bodies and loses the motive, the premise
binders and the index-equation binders; the singleton elimination
$\tscase{}{}{}{}$ collapses to its branch like $\tsplit{}{}$ and
$\tcast{}{}{}$.  Proofs are never inspected --- $\erasef$ does not
recurse into the category of proofs at all --- which is the syntactic
form of proof irrelevance.

\begin{remark}[Proofs inside index terms]
\label{rem:proof-indices}
An index term may contain proof subterms --- $\vect(f\,x\,\pi)$ for $f :
\Pi x{:}\nat.\ \varphi \to \nat$ --- and these vanish under erasure:
$\erase{f\,x\,\pi} = f\,x$.  Consequently the translated index
(\cref{def:guard-F}) and its interpretation
(\cref{def:valuation}) are independent of the proof, for \emph{arbitrary}
$f$, not only parametric ones.  This is the translation-level form of
interpreting all erased data by one canonical element, with three
precedents: Letouzey's and MetaCoq's $\bx$ (operational)
\cite{Letouzey2003,SozeauEtAl2020}, Atkey's $!_0 x = {!}_0 y$ canonical
erased realiser (algebraic) \cite{Atkey2018}, and the canonical proof
element of set-theoretic proof-irrelevant models (denotational)
\cite{LeeWerner2011}.  \Cref{lem:sem-subst}(2) is the precise statement
in our model.
\end{remark}

\begin{lemma}[Erasure commutes with substitution]
\label{lem:erase-subst}
$\erase{\subst{t}{a}{x}} = \subst{\erase{t}}{\erase{a}}{x}$;\quad
$\erase{\subst{t}{\pi}{p}} = \erase{t}$;\quad
$\erase{\subst{t}{P}{X}} = \erase{t}$.
\end{lemma}

\begin{proof}
Induction on $t$.  For the first equation the variable case is
immediate, and all other cases follow from the induction hypothesis
because $\erasef$ is homomorphic in exactly the informative positions
and substitution acts pointwise on them; positions discarded by
$\erasef$ (proofs, motives, annotations, premise and equation binders)
absorb the substitution without residue.  For the second equation,
observe that $p$ can occur in $t$ only inside proof positions or inside
types, propositions and annotations, all of which $\erasef$ discards;
similarly for the third.
\end{proof}

\begin{lemma}[Erasure simulates reduction]
\label{lem:erase-sim}
If $t \red u$ then $\erase{t} \lred \erase{u}$ or $\erase{t} =
\erase{u}$.  Consequently $t \conv u$ implies $\erase{t} \lconv
\erase{u}$.
\end{lemma}

\begin{proof}
By induction on the derivation of $t \red u$.  For root steps:
\begin{itemize}[leftmargin=2em]
\item $(\lambda x{:}\sigma.t)\,a \red \subst{t}{a}{x}$: the erasure
  $\beta$-reduces to $\erase{\subst{t}{a}{x}}$ by \cref{lem:erase-subst}.
\item $(\lambda p{:}\varphi.t)\,\pi \red \subst{t}{\pi}{p}$: both sides
  erase to $\erase{t}$; zero steps.
\item Case on a constructor:
  \begin{multline*}
    \erase{\tcase{\wcns{c}_i(\tup{a};\tup{\pi})}{\tup{z}.z.\tau}{\dots}}
    = \mathsf{case}(\wcns{c}_i(\erase{\tup{a}}); \{\wcns{c}_j(\tup{y}_j)
      \Rightarrow \erase{t_j}\}_j)\\
    \lred \subst{\erase{t_i}}{\erase{\tup{a}}}{\tup{y}_i},
  \end{multline*}
  one $\iota$-step, and by \cref{lem:erase-subst} the target equals
  $\erase{t_i[\tup{a}/\tup{y}_i, \tup{\pi}/\tup{p}_i,
  \overrightarrow{\mathsf{refl}}/\tup{q}_i]}$: the informative
  substitution is simulated pointwise and the proof substitutions (for
  $\tup{p}_i$ and $\tup{q}_i$) are invisible.
\item $\mu$ (fixpoint unfolding) and $\delta$: erase to the
  corresponding $\lbox$ steps via \cref{lem:erase-subst}.
\item The singleton rules for $\tval{}$/$\trefine{}{}$, $\tcast{}{}{}$,
  $\tsplit{}{}$ and $\tscase{J}{}{}{}$: both sides erase to $\erase{t}$
  (using \cref{lem:erase-subst} for the proof substitutions); zero
  steps.
\end{itemize}
For congruence steps: a step inside an informative position maps to a
step (or equality) in the corresponding position of the erasure by the
induction hypothesis; a step inside a proof, type, proposition, motive,
index or annotation position leaves the erasure unchanged.
\end{proof}

\begin{remark}
\Cref{lem:erase-sim} is the equational shadow of extraction correctness
\cite[\S 2.3]{Letouzey2004}; the semantic half --- the extracted term
\emph{realizes} its type --- becomes the fundamental lemma of
\cref{sec:soundness}.  Note where the indexed extension is invisible: the
$\lbox$ image of $\mathsf{vhead}$ is
$\lambda n\,v.\ \mathsf{case}(v; \{\csvnil \Rightarrow \bx \mid
\csvcons(k,a,w) \Rightarrow a\})$ --- the impossible branch has become
dead code, the possible branch a projection --- and nothing in $\lbox$
records that $v$ was drawn from an indexed family.  Soundness of the
translation cannot, therefore, be an invariant of the program half
alone; it is recovered in \cref{sec:semantics} by interpreting the
\emph{types}.
\end{remark}
